%%%%%%%%%%%%%%%%%%%%%%%%%%%%% Define Article %%%%%%%%%%%%%%%%%%%%%%%%%%%%%%%%%%
\documentclass{article}
%%%%%%%%%%%%%%%%%%%%%%%%%%%%%%%%%%%%%%%%%%%%%%%%%%%%%%%%%%%%%%%%%%%%%%%%%%%%%%%

%%%%%%%%%%%%%%%%%%%%%%%%%%%%% Using Packages %%%%%%%%%%%%%%%%%%%%%%%%%%%%%%%%%%
\usepackage{geometry}
\usepackage{url}
\usepackage{graphicx}
\usepackage{amssymb}
\usepackage{amsmath}
\usepackage{amsthm}
\usepackage{empheq}
\usepackage{mdframed}
\usepackage{booktabs}
\usepackage{lipsum}
\usepackage{graphicx}
\usepackage{color}
\usepackage{psfrag}
\usepackage{pgfplots}
\usepackage{bm}
\usepackage{hyperref} 
\usepackage[spanish]{babel}
%%%%%%%%%%%%%%%%%%%%%%%%%%%%%%%%%%%%%%%%%%%%%%%%%%%%%%%%%%%%%%%%%%%%%%%%%%%%%%%

% Other Settings

%%%%%%%%%%%%%%%%%%%%%%%%%% Page Setting %%%%%%%%%%%%%%%%%%%%%%%%%%%%%%%%%%%%%%%
\geometry{a4paper}

%%%%%%%%%%%%%%%%%%%%%%%%%% Define some useful colors %%%%%%%%%%%%%%%%%%%%%%%%%%
\definecolor{ocre}{RGB}{243,102,25}
\definecolor{mygray}{RGB}{243,243,244}
\definecolor{deepGreen}{RGB}{26,111,0}
\definecolor{shallowGreen}{RGB}{235,255,255}
\definecolor{deepBlue}{RGB}{61,124,222}
\definecolor{shallowBlue}{RGB}{235,249,255}
%%%%%%%%%%%%%%%%%%%%%%%%%%%%%%%%%%%%%%%%%%%%%%%%%%%%%%%%%%%%%%%%%%%%%%%%%%%%%%%

%%%%%%%%%%%%%%%%%%%%%%%%%% Define an orangebox command %%%%%%%%%%%%%%%%%%%%%%%%
\newcommand\orangebox[1]{\fcolorbox{ocre}{mygray}{\hspace{1em}#1\hspace{1em}}}
%%%%%%%%%%%%%%%%%%%%%%%%%%%%%%%%%%%%%%%%%%%%%%%%%%%%%%%%%%%%%%%%%%%%%%%%%%%%%%%

%%%%%%%%%%%%%%%%%%%%%%%%%%%% English Environments %%%%%%%%%%%%%%%%%%%%%%%%%%%%%
\newtheoremstyle{mytheoremstyle}{3pt}{3pt}{\normalfont}{0cm}{\rmfamily\bfseries}{}{1em}{{\color{black}\thmname{#1}~\thmnumber{#2}}\thmnote{\,--\,#3}}
\newtheoremstyle{myproblemstyle}{3pt}{3pt}{\normalfont}{0cm}{\rmfamily\bfseries}{}{1em}{{\color{black}\thmname{#1}~\thmnumber{#2}}\thmnote{\,--\,#3}}
\theoremstyle{mytheoremstyle}
\newmdtheoremenv[linewidth=1pt,backgroundcolor=shallowGreen,linecolor=deepGreen,leftmargin=0pt,innerleftmargin=20pt,innerrightmargin=20pt,]{theorem}{Theorem}[section]
\theoremstyle{mytheoremstyle}
\newmdtheoremenv[linewidth=1pt,backgroundcolor=shallowBlue,linecolor=deepBlue,leftmargin=0pt,innerleftmargin=20pt,innerrightmargin=20pt,]{definition}{Definition}[section]
\theoremstyle{myproblemstyle}
\newmdtheoremenv[linecolor=black,leftmargin=0pt,innerleftmargin=10pt,innerrightmargin=10pt,]{problem}{Problem}[section]
%%%%%%%%%%%%%%%%%%%%%%%%%%%%%%%%%%%%%%%%%%%%%%%%%%%%%%%%%%%%%%%%%%%%%%%%%%%%%%%

%%%%%%%%%%%%%%%%%%%%%%%%%%%%%%% Plotting Settings %%%%%%%%%%%%%%%%%%%%%%%%%%%%%
\usepgfplotslibrary{colorbrewer}
\pgfplotsset{width=8cm,compat=1.9}
%%%%%%%%%%%%%%%%%%%%%%%%%%%%%%%%%%%%%%%%%%%%%%%%%%%%%%%%%%%%%%%%%%%%%%%%%%%%%%%

%%%%%%%%%%%%%%%%%%%%%%%%%%%%%%% Title & Author %%%%%%%%%%%%%%%%%%%%%%%%%%%%%%%%
\title{Proyecto NHANES}
\author{Vicente Alfaro - Cristian Chávez - Isidora Guíñez - Bruno Parra}
%%%%%%%%%%%%%%%%%%%%%%%%%%%%%%%%%%%%%%%%%%%%%%%%%%%%%%%%%%%%%%%%%%%%%%%%%%%%%%%

\begin{document}
    \maketitle
    \section*{Introducción}
    \section{Pregunta de Investigación y Motivación}
\section{Fuente de Datos y Consideraciones Éticas}

\subsection{Origen de los Datos y Unidad de Observación}
Los datos utilizados en este estudio provienen de la \textit{National Health and Nutrition Examination Survey} (NHANES), un programa de estudios diseñado para evaluar el estado de salud y nutrición de adultos y niños en los Estados Unidos. Específicamente, se utiliza el ciclo de datos continuos correspondiente al período 2021-2023. La unidad de observación es una persona individual examinada \cite{cdc_nhanes_about}. 

La base de datos original se distribuye en múltiples módulos independientes en formato SAS Transport (\texttt{.xpt}). Para la consolidación de la muestra analítica, se cruzaron los registros de cuatro módulos clave utilizando el identificador único de cada paciente (\texttt{SEQN}): Demografía (\texttt{DEMO\_L}), Medidas Corporales (\texttt{BMX\_L}), Presión Arterial (\texttt{BPXO\_L}) y Colesterol Total (\texttt{TCHOL\_L}).

\subsection{Consideraciones Éticas}
El uso de esta base de datos se enmarca en un contexto estrictamente académico y estadístico. Los datos de NHANES son de dominio público y han sido previamente anonimizados por el \textit{National Center for Health Statistics} (NCHS) del CDC, garantizando que no sea posible identificar a los individuos encuestados. El análisis propuesto no busca la estigmatización de ningún subgrupo demográfico, sino comprender las relaciones entre indicadores antropométricos y cardiovasculares.

\section{EDA}

\subsection{Descripción de Variables}

% A partir de la consolidación de los módulos mencionados, a continuación se describen las variables selccionadas para el análisis:

% \begin{itemize}
%     \item \textbf{Colesterol Total (\texttt{LBXTC}):} Variable numérica continua que mide el nivel de colesterol total en la sangre, expresado en miligramos por decilitro (mg/dL). Esta medición se obtiene a partir de exámenes de laboratorio de los participantes y constituye la variable respuesta central de nuestro modelo.
% \end{itemize}


% Los predictores se dividen en tres grandes grupos fisiológicos y demográficos:

% \textbf{Predictores Demográficos:}
% \begin{itemize}
%     \item \textbf{Edad (\texttt{RIDAGEYR}):} Variable numérica que indica la edad del participante en años en el momento del examen.
%     \item \textbf{Sexo (\texttt{RIAGENDR}):} Variable categórica codificada (1 = Masculino, 2 = Femenino).
% \end{itemize}

% \textbf{Predictores Antropométricos:}
% \begin{itemize}
%     \item \textbf{Índice de Masa Corporal (\texttt{BMXBMI}):} Variable numérica continua (kg/m$^2$) calculada a partir del peso y la altura al momento de la medición.
%     \item \textbf{Circunferencia de Cintura (\texttt{BMXWAIST}):} Variable numérica continua medida en centímetros (cm), indicadora de la acumulación de grasa abdominal.
%     \item \textbf{Peso (\texttt{BMXWT}):} Variable numérica continua medida en kilogramos (kg).
% \end{itemize}

% \textbf{Predictores Clínicos:}
% \begin{itemize}
%     \item \textbf{Presión Arterial Sistólica (\texttt{BPXOSY1}):} Variable numérica continua (mm Hg) correspondiente a la primera medición de la presión sistólica (fase de contracción del corazón).
%     \item \textbf{Presión Arterial Diastólica (\texttt{BPXODI1}):} Variable numérica continua (mm Hg) correspondiente a la primera medición de la presión diastólica (fase de relajación del corazón).
% \end{itemize}
    
    \section{Plan de Modelamiento Estadístico}
    \section*{Cierre}
    \bibliographystyle{plain}
    \bibliography{referencias}
\end{document}